% ECONOMICS_PROBLEM: {"id": "E21-375", "title": "Present bias and aggregate saving under borrowing constraints", "jel_code": "E21", "source_id": "OP-0607"}
% FIRST_POSTED: 2026-10-09
% ATTEMPT_STATUS: unresolved
% ENTRY_KIND: research_attempt
\documentclass[11pt]{article}
\usepackage[T1]{fontenc}
\usepackage[utf8]{inputenc}
\usepackage[margin=25mm]{geometry}
\usepackage{amsmath,amssymb,amsthm,xurl,hyperref}
\newtheorem{proposition}{Proposition}
\newtheorem{lemma}{Lemma}
\newtheorem{corollary}{Corollary}
\newcommand{\E}{\mathbb E}
\newcommand{\R}{\mathbb R}
\newcommand{\Prb}{\mathbb P}
\newcommand{\Var}{\operatorname{Var}}
\DeclareMathOperator*{\argmax}{arg\,max}
\DeclareMathOperator*{\argmin}{arg\,min}
\setlength{\parindent}{0pt}
\setlength{\parskip}{6pt}
\title{Present bias and aggregate saving under borrowing constraints: a research attempt}
\author{GPT-6 (Codex)}
\date{9 October 2026}
\begin{document}
\maketitle
\section*{Target and outcome}
\textbf{Status: unresolved.} The catalogue seeks the stationary aggregate-asset change from removing present bias while preserving other type characteristics and recomputing equilibrium. This attempt solves a two-period overlapping-generations economy with a nonnegative saving constraint and heterogeneous present bias. It derives a unique positive stationary capital stock and an exact comparative static. This is a restricted benchmark with no competing future saving selves, not a solution for infinitely lived sophisticated households with income risk and commitment choices.

Moll's lecture \cite{moll}, whose title page dates the course to Spring 2018 despite its 2019 archive URL, explicitly motivates present bias together with borrowing constraints and portfolio choice. The derivation below is additional and does not attribute a stationary-capital theorem to the slides.

\section*{The constrained household problem}
A young household of type $(\beta,\delta,e)$ has wage income $we$, where $e>0$ is labor efficiency. It lives for two periods and chooses saving $s\geq0$. Gross return next period is $R>0$. There is no old-age endowment initially. Its young-self objective is
\[
 \log(we-s)+\beta\delta\log(Rs),\qquad 0<s<we,
\]
where $0<\beta\leq1$ and $0<\delta<1$. The old self consumes all accumulated resources; there is no later investment decision. Therefore sophisticated and naive beliefs coincide in this specific two-decision life cycle. The borrowing constraint and positive old-age consumption require nonnegative saving, and the logarithmic objective places the optimum strictly inside the stated interval.

The first-order condition and strict concavity give the unique policy
\[
 s=q(\beta,\delta)we,\qquad
 q(\beta,\delta)=\frac{\beta\delta}{1+\beta\delta}.       \tag{1}
\]
Saving is independent of $R$ because income and substitution effects cancel under logarithmic utility with zero retirement endowment. This cancellation is a special property, not a general statement about present bias.

To expose the constraint, add a promised old-age pension $z\geq0$ with fixed financing treated elsewhere in the budget. The same calculation gives
\[
 s=\max\left\{0,\frac{\beta\delta we-z/R}{1+\beta\delta}\right\}. \tag{2}
\]
At the boundary, optimality follows from the one-sided derivative. For fixed $(w,R,z,e,\delta)$, saving is weakly increasing in $\beta$, including the zero-saving region. However, a pension changes public liabilities and their financing; it cannot be inserted into the capital equilibrium below without adding those resources and claims. Equation (2) is a household comparative static, not the announced equilibrium theorem.

\section*{A closed capital-market equilibrium}
Now return to $z=0$. Each generation has mass one and an independently recurring joint type law $F$ with $\E_F e=1$. There is one unit of aggregate efficiency labor. Firms have production
\[
 Y=K^\alpha L^{1-\alpha},\qquad0<\alpha<1,
\]
full capital depreciation, and competitive factor prices. Thus
\[
 w(K)=(1-\alpha)K^\alpha,\qquad R(K)=\alpha K^{\alpha-1}.
\]
Household saving buys the productive capital used next period; there is no government debt, housing or foreign asset. Hence aggregate saving is capital by explicit market clearing. Set
\[
 Q(F)=\E_F[e\,q(\beta,\delta)].
\]
Finite $\E e=1$ and $0<q<1$ imply $0<Q<1$. Summing (1) gives the transition equation
\[
 K_{t+1}=Q(F)(1-\alpha)K_t^\alpha.                    \tag{3}
\]
A positive stationary equilibrium must therefore satisfy
\[
 K^*(F)=\{(1-\alpha)Q(F)\}^{1/(1-\alpha)}.            \tag{4}
\]
It is the unique positive solution. The degenerate zero-capital point is excluded because the maintained equilibrium has positive consumption and finite positive factor prices. Starting from any $K_0>0$, taking logarithms of (3) yields a linear recursion with coefficient $\alpha\in(0,1)$, so $K_t\to K^*$.

The stationary cross-section consists of each period's young households with zero beginning-of-period assets and old households carrying their previous saving claims. With cohort mass one, the integral of beginning assets over this two-cohort measure is $K^*$. If one normalizes total population to one instead, mean assets are $K^*/2$; comparisons must retain the same convention.

\section*{Removing present bias while retaining correlations}
The intervention changes only $\beta$ to one for each type, keeping that household's $(\delta,e)$ and their joint distribution. Consequently
\[
 Q_1=\E_F\left[e\frac{\delta}{1+\delta}\right],\qquad
 Q_0=\E_F\left[e\frac{\beta\delta}{1+\beta\delta}\right].
\]
Because the fraction is increasing in $\beta$, $Q_1\geq Q_0$, strictly if $\Pr_F(\beta<1)>0$. Recomputing wages and returns through (4) gives the exact aggregate change
\[
 \Delta A=\{(1-\alpha)Q_1\}^{1/(1-\alpha)}
          -\{(1-\alpha)Q_0\}^{1/(1-\alpha)}\geq0.      \tag{5}
\]
This is an equilibrium result inside the benchmark, rather than aggregation at fixed wages. Ignoring correlation between earnings efficiency and present bias would replace $\E(e q)$ by $\E e\E q$ and generally change its magnitude.

For homogeneous synthetic types with $\alpha=1/3$, $\delta=.9$ and $\beta=.5$, $Q_0=.45/1.45\approx.3103$, while removing present bias gives $Q_1=.9/1.9\approx.4737$. The stationary capital ratio is
\[
 K_1^*/K_0^*=(Q_1/Q_0)^{3/2}\approx1.886.
\]
This is a model illustration, not an estimate of actual missing capital. If the change occurs at date zero, the capital stock cannot jump immediately: (3) gives the transition from the common initial $K_0$, distinguishing the stationary comparison from its adjustment path.

\section*{Why the infinite-life problem remains}
An infinitely lived sophisticated agent selects consumption anticipating the policies of its future selves. Its continuation evaluation is not the value function of an ordinary discounted control problem with discount factor $\beta\delta$; $\beta$ applies once to the entire future stream. Borrowing limits, commitment accounts, income transitions and equilibrium prices can alter future-self responses. The simple proof of (5) exploits a life cycle with no later saving choice and cannot be transplanted to that setting.

The benchmark also lacks endogenous labor, uninsurable risk, bequests and asset heterogeneity. With pensions or government debt, the asset integral no longer equals productive capital without additional clearing equations. Establishing and selecting stationary equilibria in the full future-self game, and identifying the joint type law from behavior, remain unresolved. Equation (5) is therefore a rigorous special-case sign result, not a universal theorem that removing present bias always raises aggregate wealth.

\section{Completion Estimate}
42\%

This is a post hoc, heuristic estimate for the full catalogue target.
The two-period heterogeneous overlapping-generations benchmark has unique positive stationary capital and an exact bias-removal comparison preserving type correlations. Sophisticated infinite-life self-equilibria, stochastic borrowing constraints, commitment choices and empirical type identification remain outside that special case.

\begin{thebibliography}{9}
\bibitem{moll} B. Moll, Lecture 11: Good Research Topics and Open Questions, Distributional Macroeconomics, Spring 2018, slide 6. Primary PDF inspected:
\url{https://benjaminmoll.com/wp-content/uploads/2019/07/Lecture11_2149.pdf}.
\end{thebibliography}
\end{document}
